\documentclass[12 pt ,a4 paper , twoside, openright]{book}
 % use larger type; default would be 10pt
\usepackage{setspace}
\onehalfspacing
\usepackage[utf8]{inputenc} % set input encoding (not needed with XeLaTeX)
\newcommand\floor[1]{\lfloor#1\rceil}
\renewcommand\labelenumi{(\theenumi)}

%%% Examples of Article customizations
% These packages are optional, depending whether you want the features they provide.
% See the LaTeX Companion or other references for full information.

%%% PAGE DIMENSIONS
\usepackage{geometry} % to change the page dimensions
\geometry{a4paper} % or letterpaper (US) or a5paper or....
\usepackage[english]{babel}

%%% BIBLIOGRAPHY
\usepackage{natbib}
%\usepackage[babel]{csquotes}
%\usepackage[]{biblatex} % for the bibliography
%%\nocite{*}
%\usepackage[style=alphabetic]{biblatex}

% ABSTRACT ENVIRONMENT FOR BOOK CLASS
\usepackage{fancyhdr}
\newcommand{\fncyblank}{\fancyhf{}}
\newenvironment {abstract}%
{\cleardoublepage \fncyblank \null \begin{center}%
	\bfseries \abstractname \end{center}}%
{ \null }

% Page Layout

% \geometry{margin=1.5 in} % for example, change the margins to 2 inches all round
% \geometry{landscape} % set up the page for landscape
%   read geometry.pdf for detailed page layout information
\usepackage{graphicx} % support the \includegraphics command and options
%----------------------------------------------------------------------------------------------------------------------------------------------------------------------------------
%\usepackage{amsfonts}
\usepackage{amsmath}
\usepackage{amsthm}
\usepackage{amssymb}
\usepackage[parfill]{parskip} % Activate to begin paragraphs with an empty line rather than an indent

\usepackage[bookmarks]{hyperref}

\usepackage{bold-extra}

%----------------------------------------------------------------------------------------------------------------------------------------------------------------------------------

%%% PACKAGES
\usepackage{booktabs} % for much better looking tables
\usepackage{array} % for better arrays (eg matrices) in maths
\usepackage{paralist} % very flexible & customisable lists (eg. enumerate/itemize, etc.)
\usepackage{verbatim} % adds environment for commenting out blocks of text & for better verbatim
\usepackage{subfigure} % make it possible to include more than one captioned figure/table in a single float
%\usepackage{wrapfig} % make it possible to insert figure in with text around
\usepackage{caption} 
\captionsetup{tableposition=bottom,figureposition=bottom,font=small}

\usepackage{fancyref}

%%% LAYOUT
%%% HEADERS & FOOTERS
%\usepackage{fancyhdr} % This should be set AFTER setting up the page geometry
\pagestyle{fancy} % options: empty , plain , fancy
\fancyhf{}
\newcommand{\fncyfront}{% FRONT STYLE
%	\fancyhead[RO]{{\footnotesize\rightmark}}
	\fancyfoot[RO]{}
%	\fancyhead[LE]{\footnotesize{\leftmark}}
	\fancyfoot[LE]{}
	\fancyhead[RE,LO]{}
	\fancyfoot[C]{}
	\renewcommand{\footrulewidth}{0.3pt}
}

\newcommand{\fncymain}{% MAIN STYLE
	\fancyhead[RO]{{\footnotesize\MakeUppercase
			\rightmark}}
	\fancyfoot[RO]{\thepage}
	\fancyhead[LE]{{\footnotesize\MakeUppercase
			\leftmark}}
	\fancyfoot[LE]{\thepage}
	\fancyfoot[C]{}
	\renewcommand{\headrulewidth}{0.1pt}
	\renewcommand{\footrulewidth}{0.1pt}
}

\raggedbottom % no force to cover the whole page
\usepackage{todonotes}
%%% 
\usepackage{tikz,pgf} %per l'immagine di sfondo in copertina
\usetikzlibrary{matrix,positioning,decorations.pathreplacing}  % per la matrice a blocchi

%%% SECTION TITLE APPEARANCE
%\usepackage{sectsty}
%\allsectionsfont{\sffamily\mdseries\upshape} % (See the fntguide.pdf for font help)
% (This matches ConTeXt defaults)

%%% ToC (table of contents) APPEARANCE
\usepackage[nottoc,notlof,notlot]{tocbibind} % Put the bibliography in the ToC
\usepackage[titles,subfigure]{tocloft} % Alter the style of the Table of Contents
\renewcommand{\cftsecfont}{\rmfamily\mdseries\upshape}
\renewcommand{\cftsecpagefont}{\rmfamily\mdseries\upshape} % No bold!

%%%	PATH for images
\graphicspath{{./immagini/}}

%format title of section
%\usepackage{titlesec} % Used to customize the \section command
%\titleformat{\section}{\large\scshape\raggedright}{}{0em}{}[\titlerule] % Text formatting of sections
%\titlespacing{\section}{3pt}{3pt}{0.5 cm} % Spacing around sections

%-------------------------------------------------------------------
%	New Command Section
%-------------------------------------------------------------------
\numberwithin{equation}{chapter}

\theoremstyle{plain}
\newtheorem{Teorema}{Teorema}[chapter]
\newtheorem{Lemma}[Teorema]{Lemma}
\newtheorem{Proposizione}[Teorema]{Proposizione}
\newtheorem{Corollario}[Teorema]{Corollario}
%\newtheorem*{defn}{Definizione} 
\theoremstyle{definition}
\newtheorem{Osservazione}[Teorema]{Osservazione}
\newtheorem{Esempio}[Teorema]{Esempio}
\newtheorem{Esercizio}[Teorema]{Es.}
\newtheorem{Definizione}[Teorema]{Definizione}
\newtheorem{Notazione}[Teorema]{Notazione}



\setcounter{tocdepth}{1} 
%-------------------------------------------------------------------
%   PACKAGE PER IL CODICE MATLAB
%-------------------------------------------------------------------
\usepackage[framed,numbered,autolinebreaks,useliterate]{mcode}
%-------------------------------------------------------------------
%	INIZIO DEL DOCUMENTO
%-------------------------------------------------------------------
%

\begin{document}
\fncyfront
\frontmatter % SET FRONT LAYOUT STYLE


\newcommand{\titolo}
{
\bigskip
	\huge{
	Temp title
	}
}

\newcommand{\autori}
{
{
\begin{minipage}{0.4\textwidth}	
	\textbf{Group}: \\
	Francesco Finucci \\
        Stanislav Teghipco \\
        Fabio Michele De Vitis \\
        Luca Bianchi
        
\end{minipage}
\hfill
\begin{minipage}{0.4\textwidth}	
	\textbf{Supervisor}: \\
	PhD. Marco Piangerelli \\
        \\
        \\
\end{minipage}
}
}
%\hfill
%\begin{minipage}{0.4 \textwidth}
%	\textbf{Assistant Supervisor}: \\
%	Prof. Maria Letizia Corradini \\
%	Prof. Simonetta Boria
%\end{minipage}


\newcommand{\anno}
{
	\rule{15.8cm}{0.1mm}
	\centering{{\textsc{In collaboration with: Schnell S.p.A.}}}
}

\begin{titlepage}

%\begin{tikzpicture}[remember picture,overlay]
%\node [opacity=0.2,scale=1] at (current page.center) {\includegraphics{logo_unicam.png}};
%\end{tikzpicture}


\begin{center}
{{\Large{\textsc{University of Camerino}}}}
\rule[0.1cm]{15.8cm}{0.1mm}
\rule[0.5cm]{15.8cm}{0.5mm}
{\textsc{Computer Science (LM-18)}} \\
{\textsc{Complex Systems Design}}

%\vspace{0.1cm}
\titolo

\vspace{1cm}
\begin{figure*}[h]
	\centering
   	\includegraphics[width= 4.5 cm]{logo_unicam.png}
\end{figure*}

%
%\vspace{1cm}
%\begin{figure*}[h]
% \begin{minipage}[b]{7cm}
% 	\centering
%   	\includegraphics[width=6.5 cm]{logo_mmmi.png}
%	\end{minipage}
% \hspace{3mm} 
% \begin{minipage}[b]{7cm}
%  	\centering
%   	\includegraphics[width=6.5cm]{logo_mes.png}
% \end{minipage}
%\end{figure*}

\end{center}
\vfill
\autori
\vspace{0.5cm}

\anno
\end{titlepage}



%	\chapter{Introduzione}
Questa tesi presenta alcuni aspetti teorici e pratici della generazione di numeri pseudocasuali. 
In particolare, si pone l'attenzione sulla $\textit{Teoria del caos}$ e su alcuni algoritmi di generazione, i cui risultati sono stati testati tramite l'utilizzo di uno standard (NIST SP-800-22), formato da 15 test, il cui focus principale è quello di testare la validità crittografica di una sequenza pseudocasuale.

Le motivazioni che mi hanno spinto ad approfondire tale tema sono di duplice natura.
Innanzitutto, l'interesse per l'argomento è nato in aula, quando si è detto che la generazione di numeri effettivamente random è un problema irrisolto nella matematica.
Dopodiché, la curiosità mi ha spinto ad informarmi sull'argomento, arrivando così ad apprezzare sia la matematica presente dietro le quinte di un generatore pseudocasuale, sia le applicazioni pratiche: essendo io un programmatore, ho apprezzato particolarmente quelle crittografiche.

L'obiettivo di questa tesi di laurea è quello di approfondire tutto il processo di creazione di un generatore pseudocasuale. 
Si parte dalle proprietà matematiche che ci permettono di definire un generatore come tale, per poi svilupparne diversi tipi con l'obiettivo ultimo di testarne l'affidabilità. Grazie al lavoro svolto, è stato possibile studiare i comportamenti di diversi generatori in un contesto crittografico, nel quale l'affidabilità delle sequenze casuali deve essere assolutamente garantita.

La tesi è articolata in tre capitoli: 

Nel Capitolo \ref{chapter:teoriaCaos} si introducono i $\textit{sistemi dinamici}$, essenziali per definire quello che è il focus principale di questo capitolo: il $\textit{caos}$.
Per definire il caos è stato utilizzato un approccio topologico.
Il capitolo, infine, si conclude con la definizione di due sistemi dinamici, che saranno utilizzati poi per la creazione di due PRNG (Pseudo Random Numbers Generator): la $\textit{Mappa Logistica}$ ed il $\textit{Sistema di Lorenz Generalizzato}$, rispettivamente, un sistema dinamico discreto ed uno continuo.

Nel Capitolo \ref{chapter:algoGen} si applica la teoria del Capitolo \ref{chapter:teoriaCaos} per creare due generatori basati sui sistemi dinamici appena nominati, mappa logistica e sistema di Lorenz generalizzato.
Come ulteriore esercizio pratico, viene introdotto l'algoritmo di Box-Muller, che permette di generare una sequenza di numeri con distribuzione gaussiana, partendo dalle sequenze uniformemente distribuite ottenute tramite i due PRNG precedenti. 
Chiaramente, non essendo quest'ultima uniformemente distribuita, non è stato possibile applicare gli stessi test utilizzati per le altre due sequenze.
In particolare, ne vengono testate la curtosi, la simmetria e la normalità. 


Nel Capitolo \ref{chapter:chapter4} si testano i vari algoritmi.
Si parte con una definizione rigorosa di $\textit{variabili casuali}$ e $\textit{random walk}$, due oggetti che serviranno poi a comprendere i test.
Le definizioni continuano anche nella sezione successiva, in cui viene spiegato a cosa ci si riferisce quando in un test sull'ipotesi statistica si parla di ipotesi nulla, di ipotesi alternativa e di errori di tipo I e II.
Il capitolo prosegue con una spiegazione dettagliata dei 15 test utilizzati per validare i PRNG sviluppati.
Infine, si trovano le applicazioni di tali test alle sequenze generate con i due generatori sviluppati e, per completezza, anche su due funzioni fornite dal linguaggio di programmazione C++, con le varie conclusioni per ognuno dei generatori testati.

%	\input{./src/acknowledgements.tex}
	

%\listoffigures
%\listoftables

\fncymain
\mainmatter % SET MAIN LAYOUT STYLE

% CHAPTERS

\begin{abstract}
    \textit{Anomaly detection is crucial for ensuring the efficient operation and longevity of machinery across various industries. With the growth in sensor technology and data acquisition methods, machinery operations generate vast amounts of time series data which can be harnessed to monitor the health and performance of equipment. This paper presents a comprehensive approach to detect anomalies in machineries by analyzing time series data. We introduce state-of-the-art statistical and machine learning methods that cater to different types of time series structures and anomalies.}
\end{abstract}

\tableofcontents

% APPENDIXES
%\lstset{language=Matlab}
%\appendix

% BIBLIOGRAPHY

\bibliographystyle{abbrv}
\bibliography{./src/bibliography}

\end{document}